\documentclass[12pt,letterpaper]{article}
\usepackage[margin=1in]{geometry}
\usepackage[T1]{fontenc}
\usepackage{amsmath,amssymb}
\usepackage{newtxtext,newtxmath}
\usepackage{helvet}
\usepackage{enumitem}
\usepackage{caption,booktabs}
\usepackage[hidelinks]{hyperref}
\captionsetup{font=small,labelfont=bf,skip=4pt}
\renewcommand{\arraystretch}{1.15}
\setlength{\emergencystretch}{2em}
\pagestyle{empty}
\setlength{\parindent}{0pt}
\setlength{\parskip}{0pt}
\AtBeginDocument{\fontsize{12}{13.8}\selectfont}
\setlist[enumerate,1]{label=\arabic*.,leftmargin=0.25in,
labelsep=0.08in,align=left,topsep=0pt,itemsep=14pt,parsep=6pt}
\newcommand{\um}{\,\mu\mathrm{m}}
\newcommand{\uA}{\,\mu\mathrm{A}}
\newcommand{\uS}{\,\mu\mathrm{S}}
\newcommand{\uW}{\,\mu\mathrm{W}}
\newcommand{\mS}{\,\mathrm{mS}}
\newcommand{\kohm}{\,\mathrm{k}\Omega}
\newcommand{\V}{\,\mathrm{V}}
\newcommand{\MHz}{\,\mathrm{MHz}}
\newcommand{\kHz}{\,\mathrm{kHz}}
\newcommand{\pF}{\,\mathrm{pF}}
\newcommand{\Vus}{\,\mathrm{V}/\mu\mathrm{s}}

\hypersetup{pdfauthor={Anonymous contributor},pdftitle={ECEN 474/704 - Lab 8 Prelab},pdfsubject={Anonymized laboratory prelab}}
\begin{document}
{\fontsize{11}{12.65}\selectfont
\hfill October 23, 2026\par}
{\centering
{\fontsize{14}{16.1}\selectfont\sffamily\bfseries ECEN 474/704\par}
{\fontsize{14}{16.1}\selectfont Lab 8: Prelab\par}
{\fontsize{14}{16.1}\selectfont Anonymous contributor\par}
}
\vspace{12pt}
\begin{enumerate}

\item \textbf{Three-current-mirror OTA design (Fig.~8-4).}
The symmetrical OTA is required to meet Table~8-1 of the lab handout:
\begin{center}
\small
\begin{tabular}{@{}lc@{}}
\toprule
Parameter & Requirement\\
\midrule
Transconductance, $G_m$ & $>700\uS$\\
Slew rate & $>8\Vus$\\
Peak-to-peak output swing & $>1.2\V$\\
Load capacitance & $20\pF$\\
Power, excluding the bias-reference branch & $<500\uW$\\
Power supplies & $V_{DD}=0.9\V$, $V_{SS}=-0.9\V$\\
\bottomrule
\end{tabular}
\end{center}

\textbf{Solution.}

\textbf{Device parameters and model assumptions.}
The device parameters are taken from the single-finger
$W_{\mathrm{total}}/L=5/0.27\um$ rows of Table~I in the Lab~1 report:
\begin{center}
\small
\begin{tabular}{@{}lcc@{}}
\toprule
Parameter & NFET & PFET\\
\midrule
$K_P=\mu C_{ox}$ ($\mu$A/V$^2$) & 215.6 & 47.57\\
$|V_{TH}|$ (V) & 0.4330 & 0.4068\\
$\lambda$ (V$^{-1}$) & 0.1624 & 0.1945\\
$\theta$ (V$^{-1}$) & 0.9999 & 0.6558\\
\bottomrule
\end{tabular}
\end{center}
All devices have $L=0.270\um$ in the initial design. The DC
dimensions follow the square-law current model, with body effect,
mobility degradation, and DC channel-length modulation neglected.
The parameter $\lambda$ enters only through the output-resistance
estimate $r_o\approx1/(\lambda I_D)$.
The measured parameters depend on geometry and bias, so the
resulting dimensions define the initial analytical design.

All transistor widths below are \emph{total effective active widths}.
In the Lab~1 table, $W_{\mathrm{total}}$ already includes all
fingers. For equal parallel copies,
\[
W_{\mathrm{eff}}=m\,n_f\,w_{\mathrm{finger}}.
\]
The total width therefore already accounts for the finger count.
The same relation determines the physical finger and multiplicity
decomposition under the PDK width convention.

\newpage
\textbf{(a) Mirror ratio, current budget, and slew rate.}
The PMOS mirror ratio and tail current are defined as
\[
N=\frac{(W/L)_{5,6}}{(W/L)_{3,4}},\qquad I_T=I_{\mathrm{tail}}.
\]
The handout gives
\[
\mathrm{SR}\approx\frac{NI_T}{C_L},\qquad
I_{DD,\mathrm{OTA}}\approx(N+1)I_T.
\]
The bias-reference branch is excluded from the latter expression,
as in the handout. The slew-rate and power limits require
\begin{align*}
NI_T&>(20\times10^{-12})(8\times10^6)=160\uA,\\
(N+1)I_T&<\frac{500\uW}{1.8\V}=277.778\uA.
\end{align*}
Hence the feasible current interval is
\[
\frac{160}{N}\uA<I_T<\frac{277.778}{N+1}\uA.
\]
It exists only when
\[
160(N+1)<277.778N\quad\Longrightarrow\quad N>1.3585.
\]
A 1:1 PMOS mirror cannot meet both constraints in this model.
For $\boxed{N=3}$, the allowed tail-current interval is
$53.333\uA<I_T<69.444\uA$. The selected current is
\[
\boxed{I_T=I_{\mathrm{bias}}=60\uA}
\]
with a 1:1 M9--M10 bias mirror. Then
\begin{align*}
\mathrm{SR}&\approx\frac{3(60\uA)}{20\pF}=\boxed{9\Vus},\\
P_{\mathrm{OTA}}&\approx(1.8\V)(4)(60\uA)=\boxed{432\uW}.
\end{align*}
The ideal current model predicts constant OTA power within the
saturation-valid input range. Finite output conductance introduces
current variation with common-mode voltage; the DC sweep determines
the maximum power over this range. Including the separate
$60\uA$ bias-reference branch gives approximately $540\uW$.
Table~8-1 applies to the OTA power with this branch excluded.

\textbf{(b) Input transconductance and overdrive.}
At balanced input, $I_{D1}=I_{D2}=I_T/2=30\uA$ and
$G_m=Ng_{m1}$. The minimum input-device transconductance is
\[
g_{m1}>\frac{700\uS}{3}=233.333\uS.
\]
With $V_{OV1,2}=0.200\V$, the transconductances are
\[
g_{m1}=\frac{2(30\uA)}{0.200\V}=300\uS,\qquad
\boxed{G_m=900\uS>700\uS.}
\]

\newpage
\textbf{(c) Transistor dimensions and matching.}
The selected overdrives are $V_{OV}=0.200\V$ for M1--M8 and
$V_{OV}=0.150\V$ for the bias/tail devices M9 and M10.
At the balanced operating point,
\begin{align*}
I_{D1,2,3,4}&=I_T/2=30\uA,\\
I_{D5,6,7,8}&=NI_T/2=90\uA,\\
I_{D9,10}&=60\uA.
\end{align*}
The square-law current relation gives
\[
I_D\approx\frac12 K_P\frac WL V_{OV}^2,\qquad
\frac WL=\frac{2I_D}{K_PV_{OV}^2},\qquad W=L\frac WL.
\]
Using currents in $\mu$A and $K_P$ in $\mu$A/V$^2$,
\begin{align*}
(W/L)_{1,2}&=\frac{2(30)}{215.6(0.200)^2}=6.9573,\\
(W/L)_{3,4}&=\frac{2(30)}{47.57(0.200)^2}=31.5325,\\
(W/L)_{5,6}&=\frac{2(90)}{47.57(0.200)^2}=94.5974,\\
(W/L)_{7,8}&=\frac{2(90)}{215.6(0.200)^2}=20.8720,\\
(W/L)_{9,10}&=\frac{2(60)}{215.6(0.150)^2}=24.7372.
\end{align*}
\begin{center}
\small
\setlength{\tabcolsep}{5pt}
\begin{tabular}{@{}llcccc@{}}
\toprule
Devices & Type & $I_D$ ($\mu$A) & $V_{OV}$ (V) & $W$ ($\mu$m) & $L$ ($\mu$m)\\
\midrule
M1, M2 & NFET & 30 & 0.200 & 1.8785 & 0.270\\
M3, M4 & PFET & 30 & 0.200 & 8.5138 & 0.270\\
M5, M6 & PFET & 90 & 0.200 & 25.5413 & 0.270\\
M7, M8 & NFET & 90 & 0.200 & 5.6354 & 0.270\\
M9, M10 & NFET & 60 & 0.150 & 6.6790 & 0.270\\
\bottomrule
\end{tabular}
\end{center}
The displayed widths are analytical values before grid rounding.
The equal-length PMOS dimensions give
\[
\frac{W_{5,6}}{W_{3,4}}=3.
\]
A physical implementation with identical unit devices preserves
this ratio. The matched pairs are M1/M2, M3/M4, M5/M6, M7/M8,
and M9/M10. The M3--M5 and M4--M6 mirrors use consistent unit
geometry with the intended 1:3 ratio. NMOS bodies connect to
$V_{SS}$; PMOS bodies connect to $V_{DD}$.

Dummy MOS devices do not contribute to the active widths above.
NFET dummy terminals connect to $V_{SS}$ and PFET dummy terminals
to $V_{DD}$. Consistent dummy multiplicity in the schematic and
layout allows LVS to compare the same devices under the PDK
convention.

\newpage
\textbf{(d) DC bias and input common-mode range.}
The quiescent input is $V_{CM,Q}=0.300\V$ with zero differential input.
The approximate first-stage and mirror-reference voltages are
\begin{align*}
V_{S1}=V_{S2}&=0.300-(0.4330+0.200)=-0.3330\V,\\
V_{A1}=V_{A2}&=0.900-(0.4068+0.200)=0.2932\V,\\
V_B&=-0.900+0.4330+0.200=-0.2670\V,\\
V_{G9}=V_{G10}&=-0.900+0.4330+0.150=-0.3170\V.
\end{align*}
The reference $V_{o,Q}\approx V_B$ gives equal drain voltages
within the matched M5/M6 and M7/M8 pairs. Finite output
conductance determines the actual open-loop output bias; the DC
model without channel-length modulation does not uniquely fix it.
Setting $V_{o,Q}=0$ therefore requires a suitable differential
input offset or feedback arrangement.

At the reference point $V_{o,Q}=-0.2670\V$, the saturation checks are
\begin{center}
\small
\begin{tabular}{@{}lccc@{}}
\toprule
Devices & $V_{DS}$ or $V_{SD}$ (V) & $V_{OV}$ (V) & Margin (V)\\
\midrule
M1, M2 & 0.6262 & 0.200 & 0.4262\\
M3, M4 & 0.6068 & 0.200 & 0.4068\\
M5, M6 & 1.1670 & 0.200 & 0.9670\\
M7, M8 & 0.6330 & 0.200 & 0.4330\\
M9 & 0.5830 & 0.150 & 0.4330\\
M10 & 0.5670 & 0.150 & 0.4170\\
\bottomrule
\end{tabular}
\end{center}
All active devices are saturated in the assumed model.
The margin is $V_{DS}-V_{OV}$ for an NFET or $V_{SD}-V_{OV}$
for a PFET.

Keeping the input devices and tail source saturated gives
\begin{align*}
V_{CM,\min}&\approx V_{SS}+V_{OV10}+V_{TN1}+V_{OV1}\\
&=-0.900+0.150+0.4330+0.200=-0.1170\V,\\
V_{CM,\max}&\approx V_{DD}-|V_{TP3}|-V_{OV3}+V_{TN1}\\
&=0.900-0.4068-0.200+0.4330=0.7262\V.
\end{align*}
Thus the estimated interval is $-0.1170<V_{CM}<0.7262\V$,
with width $0.8432\V$. These bounds neglect body effect.
The transistor operating regions in the common-mode DC sweep
define the actual input range.

\newpage
\textbf{(e) Saturation-based output swing.}
M6 and M8 set the output-node compliance:
\begin{align*}
V_{o,\max}&\approx V_{DD}-V_{OV6}=0.900-0.200=0.700\V,\\
V_{o,\min}&\approx V_{SS}+V_{OV8}=-0.900+0.200=-0.700\V.
\end{align*}
The estimated saturation-valid interval is therefore
\[
\boxed{\Delta V_o=V_{o,\max}-V_{o,\min}=1.400\V>1.2\V.}
\]
In the handout's DC-sweep method, the endpoints occur where either
output device leaves saturation, and their difference gives the
output-swing interval. The maximum symmetric waveform also
depends on the quiescent output voltage.
For example, about $V_{o,Q}=-0.2670\V$,
\[
v_{o,\mathrm{pk}}<\min(0.4330,0.9670)\V=0.4330\V.
\]
A waveform centered near zero therefore requires bias or feedback
that places the quiescent output near zero.

\textbf{(f) Output resistance, gain, and frequency estimates.}
Using the Lab~1 $\lambda$ values at the initial length,
\begin{align*}
r_{o6}&\approx\frac{1}{0.1945(90\times10^{-6})}=57.127\kohm,\\
r_{o8}&\approx\frac{1}{0.1624(90\times10^{-6})}=68.418\kohm,\\
R_{\mathrm{out}}&\approx r_{o6}\parallel r_{o8}=31.132\kohm,\\
A_{v0}&\approx G_mR_{\mathrm{out}}=28.019\ \mathrm{V/V}
\quad(28.95\,\mathrm{dB}).
\end{align*}
With the $20$ pF load assumed to dominate output capacitance,
\begin{align*}
f_{p1}&\approx\frac{1}{2\pi R_{\mathrm{out}}C_L}
=255.61\kHz,\\
\mathrm{GBW}&\approx A_{v0}f_{p1}
=\frac{G_m}{2\pi C_L}=7.162\MHz.
\end{align*}
The factor $2\pi$ converts the handout's angular frequency to hertz.
The output-pole approximation omits internal-node poles, which
also affect the unity-gain frequency and phase margin. The
unity-feedback \texttt{iprobe}/\texttt{stb} configuration determines
the phase margin of the transistor circuit. Table~8-1 specifies no
numerical gain, GBW, or phase-margin limit.

\newpage
\textbf{(g) Design summary.}
\begin{center}
\small
\setlength{\tabcolsep}{5pt}
\begin{tabular}{@{}lcc@{}}
\toprule
Quantity & Requirement & First-order estimate\\
\midrule
Transconductance & $>700\uS$ & $900\uS$\\
Slew rate & $>8\Vus$ & $9\Vus$\\
Saturation-based output interval & $>1.2\V$ & $1.4\V$\\
OTA power, excluding bias branch & $<500\uW$ & $432\uW$\\
Load capacitance & $20\pF$ & $20\pF$\\
Supply rails & $\pm0.9\V$ & $\pm0.9\V$\\
Low-frequency voltage gain & Characterize & $28.95$ dB\\
Dominant pole & Characterize & $255.61\kHz$\\
Gain-bandwidth product & Characterize & $7.162\MHz$\\
Phase margin & Characterize & Requires \texttt{stb}\\
\bottomrule
\end{tabular}
\end{center}
The calculated transconductance, slew rate, output interval, and
OTA power satisfy the specified limits within the stated model.
Power follows the ideal current budget, while slew rate assumes
fully steered current. Mirror dynamics and parasitic capacitances
affect the actual rising and falling slopes.

The DC operating points establish the saturation-valid input
range. The maximum OTA power is evaluated over this range,
excluding the reference branch specified by the handout. The
output-device region transitions in a differential DC sweep
determine the output-swing limits.

For AC characterization, the measured transconductance is the
ratio of the output current flowing into $C_L$ to the differential
input voltage. The first-order relation is
\[
\frac{I_{C_L}(s)}{V_{id}(s)}\approx
G_m\frac{sR_{\mathrm{out}}C_L}{1+sR_{\mathrm{out}}C_L}.
\]
This current ratio tends to zero at DC, while internal poles
produce the high-frequency rolloff. The $700\uS$ specification
therefore applies to the passband maximum. An input-offset
correction establishes an active-region output bias for the
open-loop measurement. Rising and falling slew rates and the
unity-buffer phase margin complete the dynamic characterization;
parasitic extraction includes the effect of the physical layout.

The Lab~6 report provides a comparison between a first-order
calculation and transistor-level simulation: the simple amplifier's
pre-layout gain is $24.0950$ dB, compared with the calculated
$27.01$ dB. The OTA values above are analytical predictions under
the stated device-model assumptions.

\vspace{8pt}
\textbf{Sources.}
\begin{enumerate}[label={[\arabic*]},leftmargin=0.3in,
topsep=4pt,itemsep=4pt,parsep=0pt]
\item S. Palermo, \emph{ECEN 474/704 Lab 8: Operational
Transconductance Amplifiers}, Fig.~8-4, Table~8-1, and the
output-swing and stability-analysis instructions.
\item Anonymous contributor, \emph{ECEN 704 Lab 1 Report}, Table~I:
single-finger $W_{\mathrm{total}}/L=5/0.27\um$ parameter rows.
\item Anonymous contributor, \emph{ECEN 704 Lab 6 Report}, Tables~1--2.
\end{enumerate}
\end{enumerate}
\end{document}
